% =========================================================
% Proof: Compactness for Logical Consequence
% Source: volume-i/propositional-logic/notes/notes-compactness.tex
% =========================================================

\subsection*{Compactness for Logical Consequence}
\label{prf:compactness-consequence-old}

\begin{remark*}[Return]
$\leftarrow$ Back to Corollary (Compactness for Logical Consequence) in Notes
\end{remark*}

\begin{proof}
\Claim If $\Gamma \models \varphi$, then there exists a finite subset $\Gamma_0 \subseteq \Gamma$ such that $\Gamma_0 \models \varphi$.

\Given The Compactness Theorem for Propositional Logic: a set $\Gamma$ is
satisfiable if and only if every finite subset of $\Gamma$ is satisfiable.

\Goal To show: if every truth assignment satisfying $\Gamma$ also satisfies
$\varphi$, then some finite $\Gamma_0 \subseteq \Gamma$ has the same property.

\Strategy We prove the contrapositive: if no finite subset of $\Gamma$
entails $\varphi$, then $\Gamma$ does not entail $\varphi$.

Suppose that for every finite $\Gamma_0 \subseteq \Gamma$, $\Gamma_0
\not\models \varphi$. Then for each such $\Gamma_0$ there is a truth assignment
$v_0$ with $v_0 \models \Gamma_0$ but $v_0 \not\models \varphi$.

Consider the set $\Gamma' = \Gamma \cup \{\neg\varphi\}$. Every finite subset
of $\Gamma'$ has the form $\Gamma_0 \cup \{\neg\varphi\}$ for some finite
$\Gamma_0 \subseteq \Gamma$. By hypothesis, there exists $v_0$ satisfying
$\Gamma_0$ but not $\varphi$; hence $v_0 \models \neg\varphi$. So $v_0$
satisfies $\Gamma_0 \cup \{\neg\varphi\}$, making every finite subset of
$\Gamma'$ satisfiable.

By the Compactness Theorem, $\Gamma'$ is satisfiable: there is a truth
assignment $v$ with $v \models \Gamma$ and $v \models \neg\varphi$, i.e.\ $v
\not\models \varphi$. Hence $\Gamma \not\models \varphi$. \AsReq
\end{proof}

\begin{remark*}[Proof shape]
The proof is a \emph{contrapositive} reduction to the Compactness Theorem. The
key move is augmenting $\Gamma$ with $\neg\varphi$ to turn a consequence
question ($\Gamma \models \varphi$?) into a satisfiability question (is
$\Gamma \cup \{\neg\varphi\}$ satisfiable?). Compactness then applies
directly.
\end{remark*}

\begin{remark*}[Semantic content]
The corollary says that logical consequence is \emph{compact}: you can never
need infinitely many premises to derive a single conclusion. Any proof that uses
$\Gamma$ implicitly uses only finitely many formulas from $\Gamma$.
\end{remark*}

\begin{remark*}[Dependencies]
The proof depends on: the Compactness Theorem for Propositional Logic and the
definition of logical consequence ($\Gamma \models \varphi$).
\end{remark*}
